\documentclass[11pt,a4paper]{article}

\usepackage[T2A]{fontenc}
\usepackage[utf8]{inputenc}
\usepackage[russian]{babel}
\usepackage[a4paper,hmargin=1.5cm,vmargin=1.5cm]{geometry}
\usepackage{amsmath,amssymb}
\usepackage{array,booktabs}

\setlength{\parindent}{0pt}
\setlength{\parskip}{4pt}
\pagestyle{empty}

% Карточка варианта: номер, две задачи
% Размеры карточки: сетка 2x2 на всю страницу с зазором \cardgap для разрезания
\newlength{\cardgap}\setlength{\cardgap}{0.4cm}
\newlength{\cardw}\newlength{\cardh}
\newcommand{\card}[3]{%
  \fbox{\begin{minipage}[t][\cardh][t]{\cardw}
    \vspace{4pt}
    \textbf{Вариант #1}

    \smallskip
    Фамилия: \rule{5cm}{0.4pt}

    \medskip
    Найти общее решение уравнений.

    \medskip
    \textbf{1.} \ $#2$

    \medskip
    \textbf{2.} \ $#3$
  \end{minipage}}}

\begin{document}

% Страница с карточками: поля 5 мм (меньше принтер обычно не пропечатывает)
\newgeometry{margin=0.5cm}
\setlength{\cardw}{\dimexpr(\textwidth-\cardgap)/2-2\fboxsep-2\fboxrule\relax}
\setlength{\cardh}{\dimexpr(\textheight-\cardgap)/2-2\fboxsep-2\fboxrule-2pt\relax}

% ---------------------------------------------------------------------
% Страница 1: варианты для печати и разрезания
{\setlength{\parskip}{0pt}\offinterlineskip
\noindent\card{1}{y''' - 4y'' + 4y' = 0}{x^2 y'' - 2y = 0}\hfill
\card{2}{y''' + 9y' = 0}{x^2 y'' - 3xy' + 4y = 0}\par

\vspace{\cardgap}

\noindent\card{3}{y''' - 6y'' + 9y' = 0}{x^2 y'' + xy' - 9y = 0}\hfill
\card{4}{y''' + 4y'' + 5y' = 0}{x^2 y'' + 5xy' + 4y = 0}\par}

\newpage
\restoregeometry
% ---------------------------------------------------------------------
% Страница 2: ответы для преподавателя
\begin{center}
{\Large\bfseries Пятиминутка: ответы}
\end{center}

Проверяет домашнее задание к \texttt{lesson\_01.tex}: общее решение уравнения порядка выше второго и уравнение Эйлера. Время "--- 10 минут.

В каждом варианте задача 1 сводится к вынесению $\lambda$, задача 2 "--- к квадратному уравнению. Типы корней по вариантам:
\begin{center}
\renewcommand{\arraystretch}{1.3}
\begin{tabular}{@{}lll@{}}
\toprule
Вариант & Задача 1 & Задача 2\\
\midrule
1 & $0$ и кратный $2$ & простые $2$, $-1$\\
2 & $0$ и $\pm 3i$ & кратный $2$\\
3 & $0$ и кратный $3$ & простые $3$, $-3$\\
4 & $0$ и $-2 \pm i$ & кратный $-2$\\
\bottomrule
\end{tabular}
\end{center}

\subsection*{Вариант 1}

\textbf{1.} $\lambda^3 - 4\lambda^2 + 4\lambda = \lambda(\lambda - 2)^2 = 0$ \quad$\Rightarrow$\quad $\displaystyle y = C_1 + (C_2 + C_3 x)e^{2x}$

\textbf{2.} $\lambda(\lambda - 1) - 2 = \lambda^2 - \lambda - 2 = (\lambda - 2)(\lambda + 1) = 0$ \quad$\Rightarrow$\quad $\displaystyle y = C_1 x^2 + \frac{C_2}{x}$

\subsection*{Вариант 2}

\textbf{1.} $\lambda^3 + 9\lambda = \lambda(\lambda^2 + 9) = 0$, \ $\lambda = 0$, $\lambda = \pm 3i$ \quad$\Rightarrow$\quad $\displaystyle y = C_1 + C_2\cos 3x + C_3\sin 3x$

\textbf{2.} $\lambda(\lambda - 1) - 3\lambda + 4 = \lambda^2 - 4\lambda + 4 = (\lambda - 2)^2 = 0$ \quad$\Rightarrow$\quad $\displaystyle y = x^2(C_1 + C_2\ln x)$

\subsection*{Вариант 3}

\textbf{1.} $\lambda^3 - 6\lambda^2 + 9\lambda = \lambda(\lambda - 3)^2 = 0$ \quad$\Rightarrow$\quad $\displaystyle y = C_1 + (C_2 + C_3 x)e^{3x}$

\textbf{2.} $\lambda(\lambda - 1) + \lambda - 9 = \lambda^2 - 9 = (\lambda - 3)(\lambda + 3) = 0$ \quad$\Rightarrow$\quad $\displaystyle y = C_1 x^3 + \frac{C_2}{x^3}$

\subsection*{Вариант 4}

\textbf{1.} $\lambda^3 + 4\lambda^2 + 5\lambda = \lambda(\lambda^2 + 4\lambda + 5) = 0$, \ $D = 16 - 20 = -4$, \ $\lambda = -2 \pm i$ \quad$\Rightarrow$\quad $\displaystyle y = C_1 + e^{-2x}(C_2\cos x + C_3\sin x)$

\textbf{2.} $\lambda(\lambda - 1) + 5\lambda + 4 = \lambda^2 + 4\lambda + 4 = (\lambda + 2)^2 = 0$ \quad$\Rightarrow$\quad $\displaystyle y = \frac{1}{x^2}(C_1 + C_2\ln x)$

\subsection*{Типичные ошибки}
\begin{itemize}
  \item Потерян корень $\lambda = 0$, то есть слагаемое $C_1$: в ответе две константы вместо трёх.
  \item При кратном корне нет множителя $x$ (в задаче 1) или $\ln x$ (в задаче 2).
  \item В уравнении Эйлера записано $\lambda^2 + a\lambda + b$ вместо $\lambda(\lambda - 1) + a\lambda + b$.
\end{itemize}

\end{document}
